\section{Discussion and limitations}
\label{sec:discussion}

\subsection{What the present evidence supports}

The main finding is that a shared nuisance representation gives a small learned predictor a useful way to combine incomplete sensor readings. On the tested synthetic family, learning corrections to lifted Gaussian information sites improves prediction over the closed forms and ordinary transformer recipe. The width experiments show that the same sensor rule remains useful when additional sensors become available. The scalar ablation and static lifted control identify practical contributions from both the nuisance representation and the cavity inputs. These results build on products of experts, expectation propagation and learned messages~\citep{hinton2002poe,minka2001ep,heess2013learning,eslami2014just,jitkrittum2015kernel}, and connect those ideas to adaptation from a short calibration record.

The real benchmarks show the value of a flexible encoder. The small network improves on useful support fitted baselines, but the adapted transformer predicts better. Moving lifted sites to the transformer output layer recovers part of the deficit and preserves synthetic transfer. The failed noninferiority gate prevents treating that partial recovery as evidence that the readout has retained all of the transformer's real data flexibility. This distinction is substantive: the synthetic generator directly supplies the type of correlation represented by the nuisance factor, while real station and device responses can require richer interactions.

\input{generated/explore_table}

Post hoc diagnostics suggest underfitting by the small network. The transformer has lower NLL in \FTbeijingpfnbetterpct\% of Beijing PM$_{2.5}$ episodes, including differences of approximately \FTbeijinggapqQOne\ to \FTbeijinggapqQTwo\ nats in each of the three easier quartiles of closed form difficulty. Its final training NLL is also lower by \FTbeijingtraingap\ nats. Increasing nuisance dimension from one to two changes the exploratory natural mask score from \Xpliftone\ to \Xplifttwo; increasing adaptation steps fivefold improves it to \Xpliftlong, still above the transformer's \RbeijingeZerovalftpfn. These results make a limitation of the local correction inputs plausible, but they do not prove the mechanism. The readout experiment tests one response to that hypothesis, and its real outcome remains seed sensitive.

The results suggest useful next comparisons. Set processing supports interactions among variable numbers of inputs~\citep{zaheer2017deepsets,lee2019settransformer}; matching a flexible set readout and a direct covariance or precision readout would isolate the practical value of the information site construction. The 203k transformer establishes a controlled study comparison. Evaluating pretrained tabular systems~\citep{hollmann2025tabpfnv2,qu2025tabicl} would address a complementary question about current predictive performance. Pretrained TabPFN belongs to the unverified v4 study, while TabICL remains unevaluated.

\subsection{Provenance and interpretation}

The experimental record includes advance protocols, retained scores and inference replay. The five original confirmation decisions reproduce in aggregate, although factor orientation prevents exact replay of every historical score cell. The larger three seed v4 continuation also completed and its artifacts were collected. Independent verification stopped before endpoint arithmetic because required checkpoint identities had not been bound before scoring, so its five endpoints remain unverified and withheld. The quantitative findings here therefore come from v1 to v3. Appendix~\ref{app:audit} provides the full protocol outcomes, orientation diagnostics and v4 audit, including the prospective engineering repairs and their distinct version boundaries.

\subsection{Statistical and domain limits}

The synthetic family deliberately has Gaussian targets, one shared nuisance, independent idiosyncratic noise and sinusoidal response deviations. The exactness arguments concern static linear Gaussian factors, while the learned nonlinear corrections are assessed empirically. Gaussian predictions cannot represent genuinely multimodal target posteriors. Neither the ablations nor the tests establish robustness to arbitrary correlated errors, changing latent rank, informative outages or a causal missingness mechanism~\citep{rubin1976inference,little2019statistical}.

Natural missingness accounts for only \Rbeijingnatmiss\% of Beijing query readings in the PM$_{2.5}$ panel, and most stress conditions remove additional sensors artificially. The device benchmark begins with complete required measurements and then simulates dropout. Both settings are useful controlled interpolation tasks, but they give limited evidence about naturally failing networks. Pollutants from the same stations and targets from the same device are related observations, not independent domain replications. The excluded drift benchmark further narrows the evidence. No clinical, exposure decision or future forecasting benefit has been evaluated.

Paired intervals quantify task or week variation conditional on the tested fits. They do not include all uncertainty from architecture development, hyperparameter choices, training randomness, temporal dependence or repeated endpoints. The second seed analysis illustrates why a narrow task interval cannot replace independent fitted seeds, especially for the real readout result. Marginal 90\% coverage is also weaker than calibration within each station, concentration range or missingness condition. Before operational use, performance and interval behavior would need validation in the intended deployment. These limits leave a useful research finding about structured readouts, while keeping broader effectiveness claims open.
